\documentclass[11pt]{article}

\usepackage[a4paper,margin=1in]{geometry}
\usepackage{amsmath,amssymb,amsthm}
\usepackage{enumitem}
\usepackage{xcolor}
\usepackage{fancyhdr}
\usepackage{needspace}
\usepackage{physics}

\definecolor{courseblue}{RGB}{0,0,255}

\setlength{\parindent}{0pt}
\setlength{\parskip}{0.35em}
\setlength{\headheight}{14pt}

\pagestyle{fancy}
\fancyhf{}
\fancyfoot[C]{\thepage}
\renewcommand{\headrulewidth}{0pt}

\renewcommand{\ip}[2]{\left\langle #1, #2 \right\rangle}

\usepackage{bm}
\renewcommand{\vec}[1]{\bm{#1}} % Converts all \vec{} commands to bold

\DeclareMathOperator{\dist}{dist}
\DeclareMathOperator{\open}{open}
\DeclareMathOperator{\closed}{closed}
\newtheorem*{lemma}{Lemma}
\newcounter{problem}
\newenvironment{problem}{%
    \newpage
    \refstepcounter{problem}%
    \Needspace{7\baselineskip}%
    \par\vspace{0.7em}%
    \noindent\textcolor{courseblue}{\rule{\textwidth}{0.4pt}}%
    \par\vspace{0.5em}%
    \begin{list}{}{%
        \setlength{\leftmargin}{2.3em}%
        \setlength{\labelwidth}{1.8em}%
        \setlength{\labelsep}{0.5em}%
        \setlength{\topsep}{0pt}%
        \setlength{\partopsep}{0pt}%
    }%
    \item[{[\theproblem]}]
}{%
    \end{list}
    \par\vspace{0.5em}
}

\newenvironment{parts}{%
    \begin{enumerate}[label={[\alph*]},leftmargin=2.2em,itemsep=0.65em,topsep=0.65em]
}{%
    \end{enumerate}
}

% Type each solution between \begin{answer} and \end{answer}.
\newenvironment{answer}{%
    \Needspace{4\baselineskip}%
    \par\vspace{0.35em}
    \noindent\textbf{Answer:}\par
}{%
    \par\vspace{0.25em}
}

\begin{document}

\begin{center}
    {\color{courseblue}\Large\textbf{AMATH/PMATH 331: Applied Real Analysis}}
    \hfill
    {\color{courseblue}\Large\textbf{Fall 2026}}

    \vspace{1.2em}
    {\color{courseblue}\Large
    Assignment 3b Submission; due Tuesday, October 27, 2026}

    \vspace{1.2em}
    
    \textbf{Name:} Vlad Paun

    \textbf{Student number:} 21291453
\end{center}

\begin{problem}
Let $p$ be a monic polynomial of odd degree $2k+1$. This means that
\[
    p(x)=\sum_{j=0}^{2k+1}c_jx^j
\]
with $c_j\in\mathbb{R}$ for $j\in\{0,1,\ldots,2k+1\}$ and $c_{2k+1}=1$.
Let $M=\max_{0\leq j\leq 2k+1}|c_j|\geq1$.
\begin{parts}
    \item Show that $p(2M)>0$.

    \emph{Hint:} You will need to use the fact that
    $\sum_{j=0}^{2k}2^j=\frac{1-2^{2k+1}}{1-2}=2^{2k+1}-1$.
    \begin{answer}
        \[p(2M) = (2M)^{2k+1} + \sum_{j=0}^{2k}c_j(2M)^j\]
        Now, let's bound the sum term:
        \[|\sum_{j=0}^{2k}c_j(2M)^j| \leq \sum_{j=0}^{2k}|c_j|(2M)^j\]
        Since $(2M)^j > 0$ because $M \geq 1$.

        But also we have that $|c_j| \leq M$ for all $j$. So 
        \[ |\sum_{j=0}^{2k}c_j(2M)^j| \leq \sum_{j=0}^{2k}2^jM^{j+1}\]
        Now since $M \geq 1$ and $j \leq 2k$, we have $M^{j+1} \leq M^{2k+1}$. so
        \[ |\sum_{j=0}^{2k}c_j(2M)^j| \leq \sum_{j=0}^{2k}2^jM^{2k+1} = M^{2k+1}(\sum_{j=0}^{2k}2^j)\]
        \[ |\sum_{j=0}^{2k}c_j(2M)^j| \leq M^{2k+1} (2^{2k+1}-1)\]
        \[ |\sum_{j=0}^{2k}c_j(2M)^j| \leq (2M)^{2k+1} - M^{2k+1}\]
        Since $M \geq 1$, $M^{2k+1} \geq 1$:
        \[ |\sum_{j=0}^{2k}c_j(2M)^j| < (2M)^{2k+1}\]
        Therefore, 
        \[p(2M) = (2M)^{2k+1} + \sum_{j=0}^{2k}c_j(2M)^j > (2M)^{2k+1} - (2M)^{2k+1}\]
        \[p(2M) >  0\]
    \end{answer}

    \item Show that $p(-2M)<0$.
    \begin{answer}
        \begin{align*}
            p(-2M) &= (-2M)^{2k+1} + \sum_{j=0}^{2k}c_j(-2M)^j\\
            &  = \sum_{j=0}^{2k}c_j(-2M)^j  - (2M)^{2k+1}
        \end{align*}
        Now, we can reuse the same bounds as above:
        \[|\sum_{j=0}^{2k}c_j(-2M)^j | \leq \sum_{j = 0}^{2k}|c_j|(2M)^j < (2M)^{2k+1}\]
        so
        \begin{align*}
            p(-2M) &= \sum_{j=0}^{2k}c_j(-2M)^j  - (2M)^{2k+1}\\
            & < (2M)^{2k+1} - (2M)^{2k+1}\\
            p(-2M) &< 0
        \end{align*}
    \end{answer}

    \item Prove that there exists at least one point $x_o\in\mathbb{R}$
    such that $p(x_o)=0$.
    \begin{answer}
        \begin{proof}
            $p$ is a continuous map on $[-2M, 2M]$. Since $p(-2M) < 0 < p(2M)$,
             by the Intermediate Value Theorem, there exists $x_0\in(-2M,2M)\subset\mathbb{R}$ such that
             $p(x_0) = 0$.
        \end{proof}
    \end{answer}
\end{parts}
\end{problem}

\begin{problem}
Suppose that the surface of the Earth can be reasonably well-approximated
 to be a sphere $S=\{\mathbf{x}\in\mathbb{R}^3:\lVert\mathbf{x}\rVert=R\}$
for some $R>0$. It is intuitively clear, and not hard to prove, that $S$
is \emph{path-connected}. You do not have to show this.

Let $T:S\to\mathbb{R}$ be the temperature function on the surface of the
Earth. That is, $T(\mathbf{x})$ is the temperature at the point
$\mathbf{x}\in S$. Assume that $T$ is a \emph{continuous} function. Show
that there must exist at least one point $\mathbf{p}$ on the surface of
the Earth which is at exactly the same temperature at the \emph{antipodal}
point $-\mathbf{p}$. That is, prove that there exists $\mathbf{p}\in S$
such that $T(\mathbf{p})=T(-\mathbf{p})$.

\emph{Hint:} Consider the function $f:S\to\mathbb{R}$ given by
$f(\mathbf{x})=T(\mathbf{x})-T(-\mathbf{x})$.
\begin{answer}
\begin{proof}
    Consider the function $f:S\to\mathbb{R}$ given by $f(\mathbf{x})=T(\mathbf{x})-T(-\mathbf{x})$. Since $T$ is continuous, $f$ is also continuous on $S$.

    We have to show that there exists $p\in S$ such that $f(p) = 0$.

    Take arbitrary $x\in S$. Since $S$ is path connected, there exists a path between $x$ and $-x$. 
    That is, there exists continuous map:
    \[\gamma : [0,1] \to S \quad \text{ such that }\quad \gamma(0) = -x \quad \gamma(1) = x\]
    Now consider $h = f \circ \gamma : [0,1] \to\mathbb{R}$.
    Since $f$ and $\gamma$ are continuous, so is $h$.
    \[h(0) = f(\gamma(0)) = f(-x) = T(-x) - T(x) = -f(x)\]
    \[h(1) = f(\gamma(1)) = f(x)\]
    If $f(x) = 0$, we are done. Otherwise, $-f(x)$ and $f(x)$ have opposite signs.
    Therefore $0$ lies between $h(0)$ and $h(1)$.
    So by the Intermediate Value Theorem,
    \[\exists \rho\in(0,1): h(\rho) = 0\]
    Therefore, $f(\gamma(\rho)) = 0$. So let $p = \gamma(\rho)\in S$.
    \[f(p) = 0\]
\end{proof}
\end{answer}
\end{problem}

\begin{problem}
Let $f:(0,\infty)\to\mathbb{R}$ be given by $f(x)=\frac{1}{x}$.
We proved in class that $f$ is continuous on $(0,\infty)$.
\begin{parts}
    \item Show that $f$ is \emph{not} uniformly continuous on $A=(0,1)$.
    \begin{answer}

    \end{answer}

    \item Show that $f$ is uniformly continuous on $B=[1,\infty)$.
    [Note that $B$ is \emph{not compact}, so you cannot use the theorem
    from class. You need to do this directly.]
    \begin{answer}

    \end{answer}
\end{parts}
\end{problem}

\begin{problem}
Fix a subset $E\subseteq\mathbb{R}^n$. Let $f,g:E\to\mathbb{R}^m$.
\begin{parts}
    \item Suppose that $f$ and $g$ are both uniformly continuous on $E$.
    Let $\lambda,\mu\in\mathbb{R}$. Show that $\lambda f+\mu g$ is also
    uniformly continuous on $E$.

    [This shows that the space of functions from $E$ to $\mathbb{R}^m$
    which are uniformly continuous on $E$ is a linear subspace of the
    vector space of all continuous functions from $E$ to $\mathbb{R}^m$.]
    \begin{answer}

    \end{answer}

    \item Let $m=1$, so $fg:E\to\mathbb{R}$ is defined. Suppose that
    $f$ and $g$ are both uniformly continuous on $E$, and that they are
    both \emph{bounded}. Show $fg$ is uniformly continuous on $E$.

    [This shows that the space of \emph{bounded} uniformly continuous
    functions from $E$ to $\mathbb{R}$ is closed under multiplication.
    In the language of abstract algebra, it is a \emph{subring} of the
    ring of continuous functions from $E$ to $\mathbb{R}$.]
    \begin{answer}

    \end{answer}

    \item Show that the assumption that $f,g$ are bounded is necessary
    for part [b]. That is, find examples of
    $f,g:E\subseteq\mathbb{R}^n\to\mathbb{R}$ which are uniformly
    continuous on $E$, and such that at least one of $f,g$ is unbounded,
    and such that $fg$ is not uniformly continuous on $E$.
    [Take $n=1$ for simplicity.]
    \begin{answer}

    \end{answer}
\end{parts}
\end{problem}

\begin{problem}
Suppose that $f:\mathbb{R}^n\to\mathbb{R}$ is continuous on
$\mathbb{R}^n$ and that $f$ \emph{goes to zero at infinity} in the
sense of problem [5] from Assignment \#3a. Prove that $f$ is uniformly
continuous on $\mathbb{R}^n$.
\begin{answer}

\end{answer}
\end{problem}

\end{document}
